\centerline{\bf \Huge Зацикливания}

\textbf{Принцип зацикливания.} Если система может находиться лишь в конечном числе состояний, и каждое следующее состояние однозначно определяется по предыдущему, то система с некоторого момента зациклится.

\textbf{Принцип зацикливания назад.} Если в условиях принципа зацикливания каждое предыдущее состояние определяется по следующему то система зацикливается без предпериода.



\problem{1}
a) На доске отмечено несколько точек. Из каждой точки проведена ровно одна стрелка в какую-то из других точек. Докажите, что, двигаясь по стрелкам, рано или поздно начнешь ходить по циклу.

\noindent
b) Докажите, что если, кроме того, в каждую из точек ведет ровно одна стрелка, то, начав из любой точки, рано или поздно попадешь в нее снова.

\problem{2}
Жители страны Пуп Мира очень гордятся тем, что у них президентская форма правления: каждые 4 года президентом избирается либо республиканец, либо демократ. ПупМировские политологи обнаружили строгий закон, по которому определяется партийность очередного президента. Хотя этот закон засекречен "Актом о демократии", в печать просочились сведения, что партийность очередного президента полностью определяется партийностью предыдущих десяти. Докажите, что последовательность партийностей президентов зациклится, и оцените как-нибудь длину периода.

\problem{3}
Кубик Рубика выведен из первоначального состояния некоторой комбинацией поворотов. Докажите, что его можно вернуть в первоначальное состояние, выполнив эту комбинацию еще несколько раз.

\problem{4}
В стране 10 городов, некоторые из них соединены дорогой с односторонним движением, в каждый город ровно одна дорога входит и из каждого города ровно одна дорога выходит. В каждом городе находится по автомобилисту. Каждый день каждый автомобилист проезжает одну дорогу. Найдите такое минимальное $N>1$, что ровно через $N$ дней все автомобилисты будут в тех городах, из которых начинали движение, вне зависимости от того, как именно проложены дороги.

\problem{5}
В последовательности $201696 \ldots$ каждая цифра, начиная с пятой, равна последней цифре суммы четырех предшествующих цифр. Докажите, что в этой последовательности снова встретится четверка 2016.

\newpage

\hspace{3mm}

\problem{6}
Пусть $a$ и $n$ - взаимно простые числа. Докажите, что последовательность остатков от деления чисел $1, a, a^2, a^3, \ldots$ на $n$ периодическая и не содержит предпериода.

\problem{7}
В тридесятом королевстве у каждого замка и каждой развилки сходятся три дороги. Рыцарь выехал из своего замка и по очереди поворачивает то направо, то налево. Докажите, что его маршрут зациклится.

\problem{8}
Государство Элмышатия всегда существовало и всегда будет существовать. Каждый день в этом государстве либо идёт дождь, либо бушует буря, либо светит солнце. Известно, что погода в данный день однозначно определяется погодой за предшествующую четырёхдневку. Всю последнюю четырёхдневку шёл дождь. Докажите, что и до, и после этого дожливых четырёхдневок было бесконечно много.